\section{Kerangka Umum Eksperimen}
\label{sec:kerangka-umum}

Bagian ini memuat hal yang berlaku untuk kelima algoritma: struktur kode, langkah \textit{preprocessing}
yang dipakai bersama, serta protokol validasi dan reproduksibilitas. Rincian khusus setiap algoritma
dibahas pada subbabnya masing-masing.

\subsection{Struktur Kode dan Antarmuka}

Implementasi \textit{from scratch} berada di \nolinkurl{src/algoritma/}, satu folder untuk setiap
algoritma. Notebook, skrip submisi, dan dataset Kaggle berada di \nolinkurl{src/kaggle/}. Hasil
\textit{run} (model tersimpan, berkas submisi, tabel, dan gambar) ditulis ke \nolinkurl{output/}, dan
\textit{unit test} berada di \nolinkurl{test/}.

Setiap kelas mengikuti antarmuka yang mirip dengan \textit{scikit-learn}: \texttt{fit(X, y)} untuk
melatih, \texttt{predict(X)} untuk label, dan \texttt{predict\_proba(X)} untuk peluang kelas batal,
sehingga \textit{threshold} keputusan dapat diatur di luar model. Perhitungan pada implementasi
\textit{from scratch} memakai NumPy. Satu-satunya pengecualian adalah KNN, yang memakai fungsi jarak
\texttt{cdist} dari SciPy. Pustaka lain, seperti pandas dan \textit{scikit-learn}, hanya dipakai di
notebook untuk memuat data, membuat pembagian \textit{fold}, dan menjadi pembanding.

\begin{table}[H]
	\centering
	\caption{Ringkasan implementasi algoritma}
	\label{tab:ringkasan-implementasi}
	\begin{tabularx}{\textwidth}{l l l Y}
		\toprule
		\textbf{Algoritma} & \textbf{Berkas} & \textbf{Kelas} & \textbf{Format model} \\
		\midrule
		KNN & \nolinkurl{knn/knn.py} & \texttt{KNN} & \textit{Pickle} bersama \textit{preprocessor} (\nolinkurl{output/knn/model.pkl}) \\
		DTL & \nolinkurl{dtl/} (enam modul) & \texttt{DTL} & JSON (\nolinkurl{output/dtl/dtl_model.json}) \\
		\textit{Logistic Regression} & \nolinkurl{lr/lr.py} & \texttt{LogisticRegression} & NumPy \texttt{.npz} \\
		SVM & \nolinkurl{svm/svm.py} & \texttt{SVM} & \textit{Pickle} (\nolinkurl{output/svm/svm_model.pkl}) \\
		\textit{Naive Bayes} & \nolinkurl{nb/nb.py}, \nolinkurl{nb/tan.py} & \texttt{NaiveBayes}, \texttt{TANNaiveBayes} & JSON \\
		\bottomrule
	\end{tabularx}
\end{table}

DTL, \textit{Logistic Regression}, SVM, dan \textit{Naive Bayes} menyediakan metode \texttt{save} dan
\texttt{load} di kelasnya. Model KNN disimpan dengan \textit{pickle} oleh
\nolinkurl{src/kaggle/knn_submission.py}, karena isi modelnya hanya data latih dan parameter.

\subsection{\textit{Preprocessing} Bersama}

Semua model memuat data dari \nolinkurl{src/kaggle/dataset/} dan membuang kolom \texttt{id} dari
fitur. Kolom \texttt{id} hanya dipakai untuk menulis berkas submisi. Data latih berisi 27.373 baris;
setelah duplikat dibuang tersisa 27.309 baris yang dipakai untuk pelatihan dan validasi.

Setiap algoritma memiliki \textit{preprocessing} sendiri karena kebutuhannya berbeda. KNN dan SVM
membutuhkan fitur berskala sama, DTL dapat memakai atribut kategorikal dan nilai kosong secara
langsung, sedangkan \textit{Naive Bayes} membutuhkan tipe fitur (kategorikal atau Gaussian). DTL dan
SVM memakai modul fitur yang sama, yaitu \nolinkurl{src/kaggle/booking_features.py}, yang berisi
pembersihan, fitur tanggal, imputasi beserta penanda nilai kosong, fitur ``saudara'' (pemesanan lain
yang mirip pada data latih), dan standardisasi.

Pada semua algoritma, statistik \textit{preprocessing} (rata-rata, simpangan baku, modus, kategori,
\textit{target encoding}, dan sejenisnya) hanya dihitung dari data latih, lalu diterapkan ke data
validasi dan data uji. Saat validasi silang, statistik dihitung ulang di dalam setiap \textit{fold}.
Data latih dan data uji tidak pernah digabungkan.

\subsection{Protokol Validasi dan Reproduksibilitas}

Metrik utama adalah Macro F1, sama dengan metrik kompetisi (Subbab~\ref{sec:metrik}). Validasi utama
memakai \textit{stratified} 5-\textit{fold cross validation} dengan \textit{seed} 42. DTL,
\textit{Logistic Regression}, dan \textit{Naive Bayes} menjalankannya pada seluruh data latih bersih,
sedangkan KNN menjalankannya pada data \textit{development} (data latih tanpa \textit{hold-out}). SVM
memilih kernel dan $C$ lebih dulu dengan CV 3-\textit{fold} pada sampel kecil karena biaya komputasi,
lalu merinci konfigurasi akhirnya dengan CV 5-\textit{fold}. Untuk mengurangi pengaruh satu pembagian \textit{fold}, DTL dan SVM juga mengulang
validasi dengan beberapa \textit{seed} (sepuluh \textit{seed} untuk DTL dan tujuh untuk SVM).

Untuk analisis yang membutuhkan satu model tetap, misalnya perbandingan dengan \textit{scikit-learn},
analisis kesalahan, dan \textit{cost-sensitive}, KNN, \textit{Logistic Regression}, dan SVM juga
memakai pembagian \textit{hold-out} 80:20 (21.847 baris latih dan 5.462 baris validasi). Rincian
skema validasi setiap algoritma ada pada subbab \textit{modeling} dan validasinya.

Versi pustaka dipatok di \nolinkurl{requirements.txt} (antara lain NumPy 2.5.4, pandas 3.1.1, dan
\textit{scikit-learn} 1.9.1). Submisi Kaggle akhir direproduksi oleh
\nolinkurl{src/kaggle/svm_notebook.ipynb} dan \nolinkurl{src/kaggle/dtl_notebook.ipynb}
(Bagian~\ref{sec:kaggle}).
